\documentclass[11pt, a4paper]{article}
\usepackage[utf8]{inputenc}
\usepackage{geometry}
\geometry{margin=1in}
\usepackage{graphicx}
\usepackage{amsmath}
\usepackage{booktabs}
\usepackage{longtable}
\usepackage{xcolor}
\usepackage{enumitem}
\usepackage{fancyhdr}
\usepackage{titlesec}
\usepackage{hyperref}
\hypersetup{colorlinks=true, linkcolor=blue!70!black, urlcolor=blue!70!black}

% Header/Footer
\pagestyle{fancy}
\fancyhf{}
\fancyhead[L]{\textsc{TWP Logbook — Assignment 02}}
\fancyhead[R]{\textsc{CSP529}}
\fancyfoot[C]{\thepage}
\renewcommand{\headrulewidth}{0.4pt}

\titleformat{\section}{\large\bfseries}{\thesection.}{0.5em}{}[\titlerule]
\titleformat{\subsection}{\normalsize\bfseries}{\thesubsection.}{0.5em}{}

\title{
    \vspace{-1cm}
    \textbf{Writing Logbook --- Assignment 02} \\[0.3em]
    \large CSP529 Technical Writing \& Publishing \\[0.2em]
    \large \textit{Stage 2: Representations of an Operating System Process}
}
\author{
    \textbf{MT26MCS024} --- Gaurav Acharya \quad
    \textbf{MT26MCS032} --- Sandeep
}
\date{September 2026}

\begin{document}
\maketitle
\thispagestyle{fancy}

% ============================================================
\section{Assignment Overview}
% ============================================================

\begin{tabular}{@{}p{4.5cm} p{10cm}@{}}
    \toprule
    \textbf{Item} & \textbf{Detail} \\
    \midrule
    Practical Title & Stage 2: Representations of an Operating System Process \\
    Topic & Operating System Boot Process \& Resource Management \\
    Group Members & Gaurav Acharya (MT26MCS024), Sandeep (MT26MCS032) \\
    Tool Used & \LaTeX\ (with \texttt{algorithm}, \texttt{amsmath}, \texttt{booktabs}, \texttt{float} packages) \\
    Date of Completion & September 2026 \\
    \bottomrule
\end{tabular}

% ============================================================
\section{Objectives \& Learning Outcomes}
% ============================================================

The objective of this assignment was to represent a single technical idea --- the operating system boot process --- using multiple distinct representational formats in \LaTeX. The learning outcomes included:

\begin{itemize}[nosep]
    \item Writing a coherent paragraph description of a technical concept.
    \item Creating a figure representation of the idea.
    \item Constructing a structured table of system components and execution phases.
    \item Formulating a mathematical equation to model OS state transitions.
    \item Writing algorithmic pseudocode for the boot sequence.
    \item Presenting the same information as bullet points and numbered lists.
    \item Learning core \LaTeX\ constructs: \texttt{figure}, \texttt{table}, \texttt{equation}, \texttt{algorithm}, \texttt{itemize}, \texttt{enumerate}.
\end{itemize}

% ============================================================
\section{Work Performed --- Representation by Representation}
% ============================================================

\subsection{Paragraph Representation}
We composed a concise paragraph describing how an OS acts as the foundational software layer managing hardware/software, starting from the bootloader through kernel activation to user-application execution. This required distilling a complex multi-stage process into a single cohesive narrative.

\subsection{Figure Representation}
We created and included a diagram (\texttt{os\_process.png}) illustrating the layered architecture of an operating system. Key \LaTeX\ skills practised:
\begin{itemize}[nosep]
    \item \texttt{\textbackslash includegraphics} with width scaling.
    \item \texttt{figure} environment with \texttt{[H]} placement specifier (using the \texttt{float} package).
    \item \texttt{\textbackslash caption} and \texttt{\textbackslash label} for cross-referencing.
\end{itemize}

\subsection{Table Representation}
We built a four-row comparison table mapping System Components (Firmware, Bootloader, Kernel, Shell) to their Primary Functions and Execution Phases. Key \LaTeX\ skills:
\begin{itemize}[nosep]
    \item \texttt{booktabs} package for professional horizontal rules (\texttt{\textbackslash toprule}, \texttt{\textbackslash midrule}, \texttt{\textbackslash bottomrule}).
    \item Column width control using \texttt{p\{\}} specifiers.
    \item \texttt{\textbackslash renewcommand\{\textbackslash arraystretch\}\{1.4\}} for row spacing.
\end{itemize}

\subsection{Equation Representation}
We formulated a state-transition equation modelling the kernel's operation:
\begin{equation*}
    S_{t+1} = k \left( S_t, \sum_{i=1}^{h} I_i(t) + \sum_{j=1}^{s} C_j(t) \right) - \Delta_{overhead}
\end{equation*}

Each variable was then interpreted: $S_{t+1}$ (next system state), $k$ (kernel function), $S_t$ (current state), $\Sigma I_i$ (hardware interrupts), $\Sigma C_j$ (system calls), $\Delta_{overhead}$ (context-switching overhead). This exercise demonstrated that mathematical notation can concisely capture complex system behaviour that would require many paragraphs of prose.

\subsection{Algorithm Representation}
We wrote pseudocode for the OS Boot Sequence using the \texttt{algorithm} and \texttt{algpseudocode} packages:
\begin{itemize}[nosep]
    \item Conditional branching (\texttt{\textbackslash If} / \texttt{\textbackslash Else}) for POST failure handling.
    \item Sequential steps from CPU initialization through kernel activation to GUI/CLI launch.
    \item Line numbering via \texttt{\textbackslash begin\{algorithmic\}[1]}.
\end{itemize}

\subsection{Bullet Point \& Numbered List Representations}
We listed the four core OS responsibilities (Process, Memory, File System, Device Management) as bullet points, and the chronological boot sequence as a numbered list. These representations demonstrated that the same technical content can serve different purposes depending on format.

% ============================================================
\section{\LaTeX\ Skills Acquired}
% ============================================================

\begin{longtable}{@{}p{5cm} p{9.5cm}@{}}
    \toprule
    \textbf{\LaTeX\ Construct} & \textbf{What Was Learned} \\
    \midrule
    \texttt{figure} environment & Image inclusion, captioning, labelling, placement control \\
    \texttt{table} environment & Structured data presentation with professional formatting \\
    \texttt{equation} environment & Mathematical typesetting with numbering and labelling \\
    \texttt{algorithm} environment & Pseudocode with conditional logic and line numbering \\
    \texttt{itemize} / \texttt{enumerate} & Unordered and ordered list formatting \\
    \texttt{float} package & Strict placement of floats using \texttt{[H]} \\
    \texttt{booktabs} package & Professional table rules without vertical lines \\
    \bottomrule
\end{longtable}

% ============================================================
\section{Challenges Encountered}
% ============================================================

\begin{enumerate}[nosep]
    \item \textbf{Float placement:} Figures and tables initially appeared on different pages than intended. Resolved by switching from \texttt{[hbt!]} to \texttt{[H]} using the \texttt{float} package.
    \item \textbf{Equation interpretation:} Formulating a meaningful mathematical model (not just a decorative equation) required careful thinking about what the OS actually computes at each time step.
    \item \textbf{Algorithm formatting:} Getting proper indentation and conditional blocks in \texttt{algpseudocode} required understanding the nested \texttt{\textbackslash If} / \texttt{\textbackslash Else} / \texttt{\textbackslash EndIf} structure.
    \item \textbf{Escaping underscores:} Special characters in algorithm text (e.g., \texttt{OS\_Boot\_Sequence}) needed \texttt{\textbackslash\_} escaping.
\end{enumerate}

% ============================================================
\section{Key Takeaways}
% ============================================================

\begin{enumerate}[nosep]
    \item Different representational formats serve different communicative purposes: equations capture precision, algorithms describe procedure, tables organize comparisons, and prose provides context and interpretation.
    \item \LaTeX\ float management is crucial for professional document layout --- uncontrolled floats can fragment the reading flow.
    \item A technical document should not merely \emph{contain} multiple representations but should integrate them into a coherent argument where each element is introduced, referenced, and interpreted in the text.
    \item The \texttt{booktabs} package produces significantly more professional tables than default \LaTeX\ table borders.
\end{enumerate}

\end{document}
